% Section content template for individual Educative lessons
% This file contains the actual content for a specific section
% Generated from Educative API section content

% Section: Requirements for the Library Management System
% Section ID: 5635469541113856
% Section Slug: requirements-for-the-library-management-system
% Generated: 2025-12-11T15:12:12.052669

In this lesson, we will list the requirements of our library management system. This is a crucial step, as requirements define the scope of a problem. Getting them right from the interviewer and understanding them well will make the system design process smooth and easy.
\\
We'll use the notational convention to identify each requirement with a unique label ``Rn,'' where ``R'' is short for requirement and ``n'' is a natural number.

\section{Requirement collection}\label{requirement-collection}
\\
For the LMS (library management system), the requirements have been defined below:

% \vspace{-1.0\baselineskip}
\noindent
\begin{minipage}[t]{0.48\textwidth}\vspace{0pt}
\begin{itemize}
\item
\textbf{R1:} The system should store information about books and members, and maintain a complete log of all book borrowing, return, reservation, and renewal transactions.
\item
\textbf{R2:} Every book must have a unique identification number and detailed information, including rack/location in the library.
\item
\textbf{R3:} Each book should include the ISBN, title, author name, subject, and publication date.
\end{itemize}
\end{minipage}
\hfill
\begin{minipage}[t]{0.48\textwidth}\vspace{0pt}
\textit{Column component type 'MxGraphWidget' not yet supported.}
\end{minipage}

% \vspace{-1.0\baselineskip}
\noindent
\begin{minipage}[t]{0.48\textwidth}\vspace{0pt}
\begin{itemize}
\item
\textbf{R4:} A book can have multiple copies; each physical copy is a distinct book item with a unique ID.
\item
\textbf{R5:} The system must support two types of users: librarian and member, each with a library card and a unique card number.
\item
\textbf{R6:} Every user must have a library card with a unique card number.
\end{itemize}
\end{minipage}
\hfill
\begin{minipage}[t]{0.48\textwidth}\vspace{0pt}
\textit{Column component type 'MxGraphWidget' not yet supported.}
\end{minipage}

% \vspace{-1.0\baselineskip}
\noindent
\begin{minipage}[t]{0.48\textwidth}\vspace{0pt}
\begin{itemize}
\item
\textbf{R7:} Members can borrow a maximum of 10 books at a time.
\item
\textbf{R8:} A member can borrow a book for 15 days.
\item
\textbf{R9:} Only one member can reserve each book item at a time.
\item
\textbf{R10:} The system must record who issued, reserved, or renewed a book item and on which date.
\end{itemize}
\end{minipage}
\hfill
\begin{minipage}[t]{0.48\textwidth}\vspace{0pt}
\textit{Column component type 'MxGraphWidget' not yet supported.}
\end{minipage}

% \vspace{-1.0\baselineskip}
\noindent
\begin{minipage}[t]{0.48\textwidth}\vspace{0pt}
\begin{itemize}
\item
\textbf{R11:} According to policy limits, the system must allow members to renew borrowed books.
\item
\textbf{R12:} The system should notify members if a book is not returned by the due date or when a reserved book becomes available.
\item
\textbf{R13:} If a book is unavailable, a member should be able to reserve it for when it becomes available.
\item
\textbf{R14:} The system should allow users to search for books by title, author, subject, or publication date.
\end{itemize}
\end{minipage}
\hfill
\begin{minipage}[t]{0.48\textwidth}\vspace{0pt}
\textit{Column component type 'MxGraphWidget' not yet supported.}
\end{minipage}

We've identified our requirements for the problem, and in the next lesson, we will define different use cases of our library management system.

% End of section content